% =============================================================================
% TEMPORÄRE DATEI – Zitate-Übersicht zur Abstimmung mit dem Betreuer
% Kann nach Prüfung gelöscht werden.
% Erstellt: 31.05.2026 | Autor: T. Küster
% =============================================================================
%
% HINWEIS: Diese Datei ist NICHT Teil der Abschlussarbeit.
% Sie dient ausschließlich der Überprüfung der verwendeten Quellen.
%
% PROBLEM IDENTIFIZIERT:
%   \cite{idta_aas2025} wird im Text verwendet, fehlt aber in Literatur.bib!
%   --> Muss ergänzt werden, bevor die Arbeit kompiliert werden kann.
%
% =============================================================================

% --- VERWENDETE ZITATSCHLÜSSEL (alphabetisch) ---------------------------------
%
% Schlüssel              | Typ        | Kurzbeschreibung                       | Status in .bib
% -----------------------|------------|----------------------------------------|---------------
% dinspec91345           | Norm       | DIN SPEC 91345:2016 RAMI 4.0           | VORHANDEN
% iec62541               | Norm       | DIN EN IEC 62541 OPC UA (Teile 1–14)   | VORHANDEN
% idta_aas2025           | Technik.   | IDTA 01001 v3.1.2 AAS-Metamodell       | FEHLT !!!
% iso23247               | Norm       | ISO 23247-1:2021 DT-Framework Fertigung| VORHANDEN
% kagermann2013          | Bericht    | Kagermann et al., Industrie 4.0, acatech| VORHANDEN
% kritzinger2018         | Artikel    | Kritzinger et al., DT Klassifikation   | VORHANDEN
% kuhn2025aas            | Artikel    | Kuhn et al., DT und AAS, 2025          | VORHANDEN
% lechler2019            | Artikel    | Lechler et al., VIBN, Procedia CIRP    | VORHANDEN
% macenski2022           | Artikel    | Macenski et al., ROS 2, Science Robotics| VORHANDEN
% mdpi2025isaac          | Artikel    | Liu et al., DT Adaptive Robotics 2025  | VORHANDEN
% mihai2022              | Artikel    | Mihai et al., Digital Twins Survey     | VORHANDEN
% monostori2014          | Artikel    | Monostori, CPPS, Procedia CIRP 2014    | VORHANDEN
% nvidia_isaacsim        | Online     | NVIDIA Isaac Sim Dokumentation         | VORHANDEN
% nvidia_omniverse       | Online     | NVIDIA Omniverse Dokumentation         | VORHANDEN
% nvidia_physx           | Online     | NVIDIA PhysX SDK Dokumentation         | VORHANDEN
% oasis_mqtt5            | Technik.   | OASIS MQTT Version 5.0 Standard        | VORHANDEN
% pixar_openusd          | Online     | OpenUSD Core Specification             | VORHANDEN
% profanter2019          | Konferenz  | Profanter et al., OPC UA vs ROS/DDS    | VORHANDEN
% ros_urdf               | Online     | Open Robotics, URDF Dokumentation      | VORHANDEN
% stark2019              | Lexikon    | Stark & Damerau, CIRP Encyclopedia     | VORHANDEN
% vdi2222                | Richtlinie | VDI 2222 Blatt 1, Konstruktionsmethodik| VORHANDEN
% vdi3693                | Richtlinie | VDI/VDE 3693 Blatt 1, VIBN            | VORHANDEN
% vdi4499                | Richtlinie | VDI 4499 Blatt 1, Digitale Fabrik      | VORHANDEN
% vdi5200                | Richtlinie | VDI 5200 Blatt 1, Fabrikplanung        | VORHANDEN

% =============================================================================
% DETAILLIERTE EINTRÄGE (aus Literatur.bib)
% =============================================================================

% ----- 1. dinspec91345 -------------------------------------------------------
% @misc{dinspec91345,
%   author = {DIN Deutsches Institut für Normung e.V.},
%   title  = {DIN SPEC 91345:2016-04 Referenzarchitekturmodell Industrie 4.0 (RAMI4.0)},
%   year   = {2016},
%   note   = {Beuth Verlag, Berlin}
% }
% Verwendet in: Einleitung, Theoretische Grundlagen, Entwicklung, Potenziale, Fazit

% ----- 2. iec62541 -----------------------------------------------------------
% @misc{iec62541,
%   author = {International Electrotechnical Commission},
%   title  = {DIN EN IEC 62541 OPC Unified Architecture – Spezifikation (Teile 1–14)},
%   year   = {2020},
%   note   = {Geneva}
% }
% Verwendet in: Theoretische Grundlagen (OPC UA), Entwicklung (SPS-Anbindung)

% ----- 3. idta_aas2025 -------------------------------------------------------
% !!!!! FEHLT IN Literatur.bib !!!!!
% Empfohlener Eintrag:
%
% @techreport{idta_aas2025,
%   author      = {{Industrial Digital Twin Association}},
%   title       = {Specification of the Asset Administration Shell –
%                  Part 1: Metamodel, Version 3.1.2 (Release 25-01)},
%   institution = {Industrial Digital Twin Association (IDTA)},
%   number      = {IDTA 01001-3-1},
%   year        = {2025},
%   url         = {https://industrialdigitaltwin.org/content-hub/aasspecifications}
% }
%
% Verwendet in: Theoretische Grundlagen (Verwaltungsschale),
%               Potenziale (Standardkonformität), Fazit (Ausblick AAS-Integration)

% ----- 4. iso23247 -----------------------------------------------------------
% @misc{iso23247,
%   author = {International Organization for Standardization},
%   title  = {ISO 23247-1:2021 Automation systems and integration –
%             Digital twin framework for manufacturing – Part 1},
%   year   = {2021},
%   note   = {International Standard, Geneva}
% }
% Verwendet in: Theoretische Grundlagen, Potenziale

% ----- 5. kagermann2013 ------------------------------------------------------
% @techreport{kagermann2013,
%   author      = {Kagermann, Henning and Wahlster, Wolfgang and Helbig, Johannes},
%   title       = {Umsetzungsempfehlungen für das Zukunftsprojekt Industrie 4.0 –
%                  Abschlussbericht des Arbeitskreises Industrie 4.0},
%   institution = {acatech – Deutsche Akademie der Technikwissenschaften},
%   year        = {2013},
%   address     = {Frankfurt am Main}
% }
% Verwendet in: Einleitung (Ausgangssituation)
% PRÜFEN: Ist dieses Dokument noch aktuell genug? Ggf. durch neuere Quellen ergänzen.

% ----- 6. kritzinger2018 -----------------------------------------------------
% @article{kritzinger2018,
%   author  = {Kritzinger, Werner and Karner, Matthias and others},
%   title   = {Digital Twin in manufacturing: A categorical literature review},
%   journal = {IFAC-PapersOnLine},
%   volume  = {51}, number = {11}, pages = {1016–1022}, year = {2018},
%   doi     = {10.1016/j.ifacol.2018.08.474}
% }
% Verwendet in: Theoretische Grundlagen (Klassifikation DT), Entwicklung (Validierung),
%               Inbetriebnahme, Fazit
% HINWEIS: Zentrale Klassifikationsquelle – wird sehr häufig zitiert.

% ----- 7. kuhn2025aas --------------------------------------------------------
% @article{kuhn2025aas,
%   author  = {Kuhn, Thomas and Antonino, Pablo Oliveira and Bauer, Andreas},
%   title   = {Digital twin and the Asset Administration Shell},
%   journal = {Software and Systems Modeling},
%   volume  = {24}, year = {2025},
%   doi     = {10.1007/s10270-024-01255-0}
% }
% Verwendet in: Theoretische Grundlagen (AAS-Abschnitt), Fazit (Ausblick)
% PRÜFEN: Vollständige Seitenangaben vorhanden? (volume 24, aber pages fehlen im .bib)

% ----- 8. lechler2019 --------------------------------------------------------
% @article{lechler2019,
%   author  = {Lechler, Tobias and Fischer, Eva and others},
%   title   = {Virtual Commissioning – Scientific Review and Exploratory Use Cases},
%   journal = {Procedia CIRP},
%   volume  = {81}, pages = {1125–1130}, year = {2019},
%   doi     = {10.1016/j.procir.2019.03.278}
% }
% Verwendet in: Theoretische Grundlagen (VIBN), Inbetriebnahme, Potenziale, Fazit

% ----- 9. macenski2022 -------------------------------------------------------
% @article{macenski2022,
%   author  = {Macenski, Steven and Foote, Tully and others},
%   title   = {Robot Operating System 2: Design, architecture, and uses in the wild},
%   journal = {Science Robotics},
%   volume  = {7}, number = {66}, pages = {eabm6074}, year = {2022},
%   doi     = {10.1126/scirobotics.abm6074}
% }
% Verwendet in: Theoretische Grundlagen (ROS 2), Entwicklung (Dobot-Anbindung)

% ----- 10. mdpi2025isaac -----------------------------------------------------
% @article{mdpi2025isaac,
%   author  = {Liu, Y. and Sun, J. and Pan, H. and others},
%   title   = {Digital Twin-Enabled Adaptive Robotics: Leveraging Large Language
%              Models in Isaac Sim for Unstructured Environments},
%   journal = {Machines},
%   volume  = {13}, number = {7}, pages = {620}, year = {2025},
%   doi     = {10.3390/machines13070620}
% }
% Verwendet in: Theoretische Grundlagen (Isaac Sim), Potenziale, Fazit (Ausblick)
% PRÜFEN: Autoren vollständig? ("others" ist Platzhalter – vollständige Autorenliste ergänzen)

% ----- 11. mihai2022 ---------------------------------------------------------
% @article{mihai2022,
%   author  = {Mihai, Stefan and Yaqoob, Mahnoor and others},
%   title   = {Digital Twins: A Survey on Enabling Technologies, Challenges, Trends},
%   journal = {IEEE Communications Surveys & Tutorials},
%   volume  = {24}, number = {4}, pages = {2255–2291}, year = {2022},
%   doi     = {10.1109/COMST.2022.3208773}
% }
% Verwendet in: Einleitung (Problemstellung), Theoretische Grundlagen (Defizite)

% ----- 12. monostori2014 -----------------------------------------------------
% @article{monostori2014,
%   author  = {Monostori, László},
%   title   = {Cyber-physical production systems: Roots, expectations and R&D challenges},
%   journal = {Procedia CIRP},
%   volume  = {17}, pages = {9–13}, year = {2014},
%   doi     = {10.1016/j.procir.2014.03.115}
% }
% Verwendet in: Einleitung (CPPS-Definition), Theoretische Grundlagen

% ----- 13. nvidia_isaacsim ---------------------------------------------------
% @online{nvidia_isaacsim,
%   author  = {NVIDIA Corporation},
%   title   = {NVIDIA Isaac Sim Documentation},
%   year    = {2025},
%   url     = {https://docs.isaacsim.omniverse.nvidia.com/5.1.0/index.html},
%   urldate = {2026-05-26}
% }
% Verwendet in: Einleitung, Theoretische Grundlagen, Methodik, Entwicklung

% ----- 14. nvidia_omniverse --------------------------------------------------
% @online{nvidia_omniverse,
%   author  = {NVIDIA Corporation},
%   title   = {NVIDIA Omniverse Platform Documentation},
%   year    = {2025},
%   url     = {https://docs.omniverse.nvidia.com/},
%   urldate = {2026-05-26}
% }
% Verwendet in: Einleitung, Theoretische Grundlagen, Inbetriebnahme, Potenziale

% ----- 15. nvidia_physx ------------------------------------------------------
% @online{nvidia_physx,
%   author  = {NVIDIA Corporation},
%   title   = {NVIDIA PhysX SDK Documentation},
%   year    = {2025},
%   url     = {https://nvidia-omniverse.github.io/PhysX/physx/latest/},
%   urldate = {2026-05-26}
% }
% Verwendet in: Theoretische Grundlagen (PhysX), Entwicklung (Kinematik)

% ----- 16. oasis_mqtt5 -------------------------------------------------------
% @techreport{oasis_mqtt5,
%   author      = {OASIS},
%   title       = {MQTT Version 5.0 – OASIS Standard},
%   institution = {OASIS},
%   year        = {2019},
%   url         = {https://docs.oasis-open.org/mqtt/mqtt/v5.0/mqtt-v5.0.html}
% }
% Verwendet in: Theoretische Grundlagen (MQTT)

% ----- 17. pixar_openusd -----------------------------------------------------
% @online{pixar_openusd,
%   author  = {Pixar Animation Studios and Alliance for OpenUSD},
%   title   = {Universal Scene Description (OpenUSD) – Core Specification},
%   year    = {2024},
%   url     = {https://openusd.org/release/intro.html},
%   urldate = {2026-05-26}
% }
% Verwendet in: Theoretische Grundlagen (OpenUSD), Entwicklung, Inbetriebnahme, Potenziale

% ----- 18. profanter2019 -----------------------------------------------------
% @inproceedings{profanter2019,
%   author    = {Profanter, Stefan and Tekat, Ayhun and others},
%   title     = {OPC UA versus ROS, DDS, and MQTT: Performance Evaluation},
%   booktitle = {Proc. IEEE ICIT 2019},
%   pages     = {955–962}, year = {2019},
%   doi       = {10.1109/ICIT.2019.8755050}
% }
% Verwendet in: Theoretische Grundlagen (Protokollvergleich), Inbetriebnahme (Diskussion)

% ----- 19. ros_urdf ----------------------------------------------------------
% @online{ros_urdf,
%   author  = {Open Robotics},
%   title   = {URDF – Unified Robot Description Format},
%   year    = {2024},
%   url     = {https://wiki.ros.org/urdf},
%   urldate = {2026-05-26}
% }
% Verwendet in: Theoretische Grundlagen, Methodik, Entwicklung (Dobot URDF), Anhang

% ----- 20. stark2019 ---------------------------------------------------------
% @incollection{stark2019,
%   author    = {Stark, Rainer and Damerau, Thomas},
%   title     = {Digital Twin},
%   booktitle = {CIRP Encyclopedia of Production Engineering},
%   publisher = {Springer, Berlin Heidelberg}, year = {2019},
%   doi       = {10.1007/978-3-642-35950-7_16870-1}
% }
% Verwendet in: Theoretische Grundlagen (Definition DT nach Stark/Damerau)

% ----- 21. vdi2222 -----------------------------------------------------------
% @misc{vdi2222,
%   author = {Verein Deutscher Ingenieure},
%   title  = {VDI 2222 Blatt 1 Konstruktionsmethodik – Methodisches Entwickeln},
%   year   = {1997}, address = {Düsseldorf}
% }
% Verwendet in: Methodik
% PRÜFEN: Gibt es eine neuere Ausgabe? Die aktuelle Fassung könnte neuer sein.

% ----- 22. vdi3693 -----------------------------------------------------------
% @misc{vdi3693,
%   author = {VDI und VDE},
%   title  = {VDI/VDE 3693 Blatt 1 Virtuelle Inbetriebnahme – Modellarten und Glossar},
%   year   = {2016}, address = {Düsseldorf}
% }
% Verwendet in: Theoretische Grundlagen (VIBN), Inbetriebnahme, Potenziale, Fazit

% ----- 23. vdi4499 -----------------------------------------------------------
% @misc{vdi4499,
%   author = {Verein Deutscher Ingenieure},
%   title  = {VDI 4499 Blatt 1 Digitale Fabrik – Grundlagen},
%   year   = {2008}, address = {Düsseldorf}
% }
% Verwendet in: Theoretische Grundlagen (Digitale Fabrik), Potenziale (Ökonomie)
% PRÜFEN: Gibt es eine neuere Ausgabe seit 2008?

% ----- 24. vdi5200 -----------------------------------------------------------
% @misc{vdi5200,
%   author = {Verein Deutscher Ingenieure},
%   title  = {VDI 5200 Blatt 1 Fabrikplanung – Planungsvorgehen},
%   year   = {2011}, address = {Düsseldorf}
% }
% Verwendet in: Theoretische Grundlagen (Fabrikplanung), Fazit

% =============================================================================
% HANDLUNGSBEDARF – ZUSAMMENFASSUNG
% =============================================================================
%
% KRITISCH (Arbeit kompiliert nicht):
%   [1] idta_aas2025 fehlt in Literatur.bib
%       Vorgeschlagener Eintrag s.o. (Abschnitt 3)
%
% ZU PRÜFEN (inhaltlich/formal):
%   [2] kagermann2013    – Noch ausreichend aktuell? Ggf. neuere Quelle ergänzen.
%   [3] kuhn2025aas      – Seitenangaben im .bib-Eintrag fehlen (pages = {...})
%   [4] mdpi2025isaac    – "others" als Platzhalter – vollständige Autorenliste?
%   [5] mihai2022        – "others" als Platzhalter – vollständige Autorenliste?
%   [6] profanter2019    – "others" als Platzhalter – vollständige Autorenliste?
%   [7] vdi2222          – Jahrgang 1997 – gibt es eine aktuellere Fassung?
%   [8] vdi4499          – Jahrgang 2008 – gibt es eine aktuellere Fassung?
%
% EMPFOHLEN (aus den eigenen Projektarbeiten des Verfassers):
%   [9]  Eigene Projektarbeit: Dobot-Tutorial (Küster, 2026) – noch nicht zitiert
%   [10] Eigene Projektarbeit: Omniverse-Tutorial / Maze Runner (Küster, 2026)
%        – noch nicht zitiert, aber inhaltlich Grundlage von Kap. 4 und 5
%        Empfehlung: Als unveröffentlichter Bericht (unpublished) oder
%        interner Laborbericht zitieren, z.B.:
%
%        @unpublished{kuester2026dobot,
%          author = {Küster, Tobias},
%          title  = {Anleitung zur Erstellung eines Digitalen Zwillings –
%                    Dobot Magician an NVIDIA Isaac Sim},
%          note   = {Projektarbeit, Hochschule Hannover, LFP},
%          year   = {2026}
%        }
%
%        @unpublished{kuester2026mazerunner,
%          author = {Küster, Tobias},
%          title  = {Anleitung zur Erstellung eines Digitalen Zwillings –
%                    Maze Runner Anlage mit NVIDIA Omniverse},
%          note   = {Projektarbeit, Hochschule Hannover, LFP},
%          year   = {2026}
%        }
%
%   [11] Quelle des Maze-Runner-Labors:
%        Labor Mechatronik (Trainingslevel Konstruktion und Inbetriebnahme
%        einer Maze Runner Produktionslinie, SS2024, HsH)
%        – bisher nicht im .bib, sollte ergänzt werden:
%
%        @misc{hsh_mazerunner2024,
%          author       = {{Hochschule Hannover, Labor Mechatronik}},
%          title        = {Trainingslevel: Konstruktion und Inbetriebnahme einer
%                          Maze Runner Produktionslinie unter Einsatz
%                          interoperabler Edge- und Cloud-Technologien},
%          howpublished = {Laborunterlagen},
%          year         = {2024},
%          note         = {Sommersemester 2024}
%        }
%
% =============================================================================
% ENDE DER DATEI
% =============================================================================
